\section{Developer handbook}
\subsection{Architecture and boundaries}
The engine owns pure rules: TOML parsing, dependency/level eligibility, streaks,
and validation. Storage owns SQLite migrations and user-scoped access. The
\texttt{Tracker} application layer implements shared use cases. CLI and FastAPI
are transports over those use cases; the vanilla JavaScript PWA consumes the API.
Subject content belongs in packs, not engine conditionals. Verifier adapters own
external effects and return bounded, sanitized evidence results.

Key files: \texttt{ascend/engine/packs.py} defines immutable Pack, Level and Node
models; \texttt{levels.py} computes candidates; \texttt{streak.py} derives streaks;
\texttt{ascend/application.py} coordinates use cases;
\texttt{ascend/storage/\_\_init\_\_.py} scopes storage;
\texttt{ascend/server/app.py} declares HTTP routes and security headers;
\texttt{ascend/cli/remote.py} implements standard-library remote access;
\texttt{web/app.js} and \texttt{web/style.css} own interaction and presentation.

\subsection{State and data flow}
A pack install parses and validates before replacing the private text file and
recording its content fingerprint. Stable IDs join task state to pack data.
Loading a changed file revokes stale review/trust. User variable settings are
applied separately from authored placeholders. A task split creates application
nodes while keeping the pack file unchanged. Prerequisites and main-level gates
control recommendations; explicit overrides are logged as overrides.

Logging resolves an account/mission/node, validates learning text, optionally runs
a verifier, then records the session in a transaction. Failed completion still
records partial work. Successful completion records its prior state and revision
so undo can avoid clobbering a later change. Streaks and time statistics are derived
from entries, not independent mutable counters. File/Git evidence timestamps can
be backdated to actual work time and are deduplicated by the existing evidence flow.

\subsection{Challenge preview and started workspace}
\path{ascend/server/presentation.py} removes instructions, downloads, hints,
completion criteria and verifier internals from web overview and
listing responses, including the current level nested in the dashboard response.
Resource links, labels and notes remain available before Start so learners can
prepare without starting the timer. \texttt{preview} and \texttt{requirements} are authored Node fields with the usual
variable and series substitutions. Older packs can omit them and show a title.
The challenge route returns full content only when the current user has a task
session or historical completion. A direct URL cannot skip the preview; an ended
session retains review access. Public pack exports and the standalone CLI remain
authoring/practice tools, so this is not an anti-cheat secrecy boundary.

In \path{web/app.js}, \texttt{showNode} builds either the short preview or the full
workspace. Start creates the existing server-clock session and navigates to the
challenge route; opening or reviewing a preview does not create a session. Reload
and Review attempt reuse existing state. No pause or timer reset is introduced.
API regression tests exercise nested/list responses and review access; browser
tests cover direct links, downloads, copyable commands and timer restoration at
mobile and desktop widths.

\subsection{Security and operational assumptions}
Every business query is user-scoped. Passwords use scrypt and server bearer tokens
are hashed. CLI token storage uses restrictive permissions; browser tokens are
session-scoped unless the learner chooses Remember me. The PWA uses DOM text construction for authored content. The server
sets a self-only content security policy, same-origin write checks and no-store
headers for API responses. Signup is opt-in in multi-user mode and has a bounded
per-process creation limit. One-use recovery codes rotate atomically and revoke
earlier tokens. Public-production limits remain those documented for the chosen
deployment topology. The public catalog never reads private installed packs.

Command templates are split before substitution and run without a shell. Plugins
run in a separate process with a timeout, but retain the host user's permissions.
Multi-user executable/network checks are restricted unless the administrator opts
in. Do not infer sandboxing from process isolation or from pack validation. The
verifier chapter gives the exact extension contract.

\subsection{Build, test and contribution loop}
Install Python 3.11 or newer and its virtual-environment support. Documentation
builds use XeLaTeX, the extra LaTeX packages and DejaVu fonts. On Debian/Ubuntu,
the documentation prerequisites can be installed with:
\begin{verbatim}
sudo apt-get install texlive-xetex texlive-latex-extra fonts-dejavu
\end{verbatim}
Pandoc was used for the one-time migration of historical Markdown; it is not
needed for routine builds or edits to the canonical LaTeX chapters.
\begin{verbatim}
python3 -m venv .venv
.venv/bin/pip install -e '.[server,dev,push]' setuptools wheel
.venv/bin/playwright install chromium
.venv/bin/pytest -v
.venv/bin/ruff format --check ascend tests scripts setup.py
.venv/bin/ruff check --select F401,F841 ascend tests scripts setup.py
.venv/bin/python scripts/build_registry.py
.venv/bin/python scripts/build_docs.py
\end{verbatim}
Tests use temporary profiles and local test servers. They cover pure rules,
application semantics, CLI/remote behavior, account isolation, migrations,
notifications, browser flows and built-wheel assets. Add meaningful regression
coverage for a behavioral change; do not run external lab attacks in the test
suite. Resource link checks are explicit network operations, separate from offline
unit tests. The CI matrix covers Python 3.11--3.13.

With an already built local image, run the optional operations drill:
\begin{verbatim}
.venv/bin/python scripts/restore_smoke.py ascend:local
\end{verbatim}
It creates two disposable Compose projects, backs up a synthetic profile, restores
it at the same internal path, runs the guide's exact reminder-disable program,
checks data and HTTP health, then removes only its own test volumes. It does not
send notifications or touch the production profile. The retained result names
the exact image tested; it is not an arm64 or external delivery check.

The \texttt{setup.py} build hook copies canonical \texttt{web/} and
\texttt{packs/community/} assets into wheels. Verify installed artifacts as well
as editable installs. \texttt{scripts/screenshots.py} uses disposable illustrative
data, not the owner's profile. Never commit databases, credentials, environment
files, personal logs, or real deployment addresses. Make feature-sized commits
and update the implementation report at each milestone.

\subsection{Maintaining the curriculum and documentation}
Canonical runtime learning content is TOML under \texttt{packs/community/}.
Files are grouped into \texttt{foundations/}, \texttt{offensive/},
\texttt{defensive/}, \texttt{engineering/}, \texttt{specialist/},
\texttt{practice/}, and \texttt{projects/}. The registry records each relative
path so installation by pack ID remains independent of the folder layout.
The exported cyber pack files must match it. Preserve IDs, increment pack version,
regenerate registry checksums, validate the dependency graph, and document revised
completion requirements. The assessment map in
\texttt{docs/curriculum/assessment-map.json} records one explicit mechanism,
experiment, transfer case, control and delayed-retrieval prompt for each main
level. Source evidence, access dates and limits are in
\texttt{docs/curriculum/sources.json}. Keep both synchronized with the authored
milestones; the curriculum regression suite checks their structural contract.

The canonical narrative documentation is LaTeX under \texttt{docs/latex/}.
\path{scripts/build_docs.py} refreshes generated command/schema/curriculum
appendices and compiles the report, complete handbook and focused guides.
Generated sections are marked as such. Edit narrative chapters directly; do not
regenerate them from legacy Markdown. The remaining Markdown documents are
convenience copies or historical inputs. Command documentation may be duplicated
there, but the LaTeX reference is complete. Historical designs are separately
compiled and explicitly labeled as superseded, never mixed with current guidance.

\subsection{Schema and extension contracts}
The pack-format and verifier chapters specify all authored fields and extension
boundaries. The API chapter combines a human-readable endpoint reference with
request-model information extracted from OpenAPI. The generated storage appendix
reproduces the numbered migration sources used by the running code. Migrations
are sequential transactions over a WAL-backed database; do not open a newer
schema with an older binary. Change shared semantics in the application layer and
check both standalone and remote callers before release.


\subsection{Linux local labs and bounded grading}
\path{ascend/local_labs/catalog.py} is the separate reference and recipe catalog
shared by tracks, the Linux labs tab and CLI. The research register is published
with \path{scripts/build_lab_catalog.py}; packaged runtime data lives in
\path{ascend/local_labs/resources}. Adding a reference never activates an adapter.
Recipe activation requires immutable downloads, reviewed code, versioned checks,
platform measurements and lifecycle evidence. The HTTP command gateway rejects
client-only runtime operations, even on a single-user server.

Migration 0011 adds account/nonce/version-bound external lab attempts and the
foreground-session guards. Migration 0012 adds assigned policy cases. Local run
state, download cache and saved work are separate and bound to the local profile;
the API cannot launch the companion or read a submitted local path. The Docker
adapters inventory resources before provisioning, verify ownership labels,
enforce internal target networks and loopback proxies, retain failed cleanup
state and reconcile after process interruption. Source/image eviction is scoped
to acquired, unused objects; user work is retained.

\path{ascend/local_labs/prefix_worker.py} is a fixed, data-only IPv4 policy
checker invoked with isolated Python startup in a temporary directory. It
accepts at most 64 KiB and 32 declarative rules, with Linux CPU, address-space,
process, descriptor, output and wall-clock limits. The parent checks its pinned
hash, output version and complete checkset and reaps failures. Partitioning route
indices at every rule boundary proves equivalence over the supported policy
language; this is not a general sandbox for executable submissions. Online cases
are server-issued; offline case keys derive a canonical case on the server.
Source, recipe, checker, worker hash and case bind the fixture identity.

\path{ascend/local_labs/upstream.py} adapts selected Labtainers/netlab report
semantics and strict specialist observation templates without launching those
providers. Empty goals and unsupported/warning exits never become passing
grades. Manual observations stay explicitly manual, including claimed environment
hashes and model trial counts. Temporary server policy tests and local reports
issue no track proof or reward. Consult
\path{scripts/verify_local_file_lab.py} and \path{scripts/verify_hosted_labs.py}
for reproducible runtime checks, and \path{docs/design/shared-lab-runner-evaluation.md} for the remaining capability
gates and Chall-Manager evaluation.

\path{ascend/local_labs/bgp.py} translates three selected BGP Labs fixtures into
bounded, non-root FRR control-plane containers. The immutable image and source
license are independently pinned. All capabilities are dropped, roots are read
only, and the internal network publishes no ports. Named volumes retain router
configuration; authored negative probes restore external controls afterward.
The adapter and launcher hashes bind the recipe's fixture identity.
Route checks require a stable customer/withheld-prefix control pair after
session establishment. Protection checks reload the saved learner configuration
before probing so earlier TTL settings cannot determine fresh-socket behavior.
The launcher uses the original absolute expiry, including after stop/resume.

\path{ascend/local_labs/workspaces.py} prepares learner-managed lab workspaces
without reserving a foreground session or launching infrastructure.
\path{ascend/local_labs/collectors.py} provides bounded read-only tool adapters
with explicit provider/device/account selection and secret-value redaction.
\path{ascend/local_labs/evidence.py} checks eleven artifact contracts, including
paired identities, distinct objects, negative controls and repeated model trials.
The server repeats the same data checks and never executes uploaded code.
Artifact consistency is distinct from provenance; these receipts retain local
authority, zero XP and no verified-track equivalence. Finished evidence workspaces
are immutable, reset creates a separate workspace and synchronization is explicit.

\subsection{Independent activity architecture}
\path{ascend/playground/catalog.py} contains versioned definitions with stable
reward identities, roles, skill prerequisites, preparation references, estimates,
device/setup needs, formats and help policies. \path{ascend/playground/fixtures.py}
keeps public evidence separate from private expected answers and debriefs. It
reuses native access and subnet generation without calling track marking. Existing
installed pack copies and assessment receipts are never rewritten.

\path{ascend/engine/practice.py} implements pure readiness, review intervals of
1/3/7/14 days, and recommendation ranking. Input is explicit skill evidence and
attempt history; total XP and mission progress are not readiness inputs. Examples
introduce skills; qualifying independent cases and transfer checks support
readiness. A setup report is saved work, not a failed conceptual check.

Migration 0008 adds user-scoped activity runs, evidence, bookmarks and reward
records, plus an optional activity-run reference on entries. Runs retain the full
versioned definition, fixture, meaningful fixture identity, prerequisite snapshot,
work, checks, hints and reveal state. The public API uses explicit projections and
never sends the private rubric before completion or authorized review. Authored
variant exhaustion is familiar practice. Semantic fixture identities prevent
cosmetic request renaming from creating new access cases.

Migration 0009 adds accumulated elapsed time, active clock intervals, per-skill
tour preferences, optional session feedback and local selection/resume/return
observations. Migration files are captured when the storage module loads, so a
running old server cannot apply a newly written schema ahead of its matching code.
Restart the server after updating backend code. The browser explains stale-server
responses instead of presenting raw missing-field errors.

\path{ascend/playground/guidance.py} supplies authored steps only for examples.
Tour writes never affect completion, rewards or readiness; autosaves preserve the
server's current tour position. \path{ascend/playground/insights.py} keeps user-scoped
feedback and a 14-day summary. Personal estimates require three distinct completed,
clear, problem-free reports and use their 75th percentile with the authored estimate
as a floor. Latest misses, due review, unwanted repeats and exhausted cases inform
ranking. Examples cannot delay assessment review or establish independent mastery.

Migration 0010 stores authored checkpoint completion separately, referencing its
partial learning entry and retaining a definition snapshot. Undoing that entry
removes its checkpoint via a foreign key. \path{ascend/checkpoints.py} requires an
open parent and writes an ordinary partial log; it never calls marking or timing.
Checkpoint durations describe parts, while Today retains the parent's full estimate.

Starts, completion, evidence and reward writes use SQLite immediate transactions.
Partial indexes and cross-table triggers enforce one foreground activity or track
session. Starting Playground requires an explicit scored-session end when needed;
starting a track requires saving the active activity. Completion is idempotent;
ledger uniqueness enforces stable first-solve and family/day practice awards.
Activity logs write a null mission directly rather than using the active-mission
fallback in ordinary logging. A unique run/day entry suppresses duplicate automatic
and manual habit events. Saved work may use last-write-wins edits; final results and
awards remain immutable. Revealed/familiar evidence cannot demonstrate a fresh
independent solve.

The shared \texttt{Tracker.playground} service supplies local CLI, remote command
transport and dedicated HTTP routes. \path{web/playground.js} renders catalog,
editable workspaces, preparation, limited hints, review and result actions using
text DOM construction. Browser drafts save to the server and use session storage
only for unsaved fallback. The service worker includes the new script; the normal
wheel builder copies canonical web assets and includes the Python content package.

Activity URLs use \texttt{\#activity/RUN}; legacy Playground run links redirect.
The focused page omits catalog tabs and finished track-session chrome. Examples
have an inline keyboard-accessible walkthrough and a specific reflection prompt;
challenges keep independent evidence questions. Assessed challenge elapsed time
is anchored to the original creation time and continues while away, until
completion or explicit end. Save cannot pause it. Examples and free workspaces
retain pause/resume persistence. Existing fixed activity rewards still apply.
History groups attempts by activity ID across versions, with individual review
at \texttt{\#attempts/ACTIVITY}; raw API history retains every run.

Run \texttt{pytest tests/test\_playground.py} for migration, ownership, reward,
concurrency, readiness, privacy and generated-fixture checks. Run
\path{scripts/playground_smoke.py} with Playwright for the desktop/phone flow and
screenshots. Broader subjects, provisioned execution environments and user pilots
for timing/readiness calibration remain follow-up work.

\subsection{Account-scoped browser runner}
\path{ascend/local_labs/hosted.py} controls reviewed adapters in separate
account-scoped profiles. \path{ascend/local_labs/runner.py} supplies the private
HTTP service, bounded request schema, machine authentication and restart sweep.
\path{ascend/server/lab_runner.py} derives account scopes from the installation
identity and authenticated account, then forwards only fixed control operations.
The web process never creates the provider runtime and continues to reject local
CLI operations. There is no arbitrary shell, submitted executable or target URL.
A single worker records queued operations and outcomes; process ownership is
locked for each runner home. Inventory precedes provisioning. Failed exports
retain compute; failed deletions retain owned resources and capacity reservations.
The 30-second reconciliation loop and container deadlines cover expired runs.

\path{web/local-labs.js} presents compact Playground cards and a shareable
\texttt{\#lab/ID} detail route. Recipe links, browser history and catalog filter
retention are tested on desktop and phone. Details contain setup, commands,
evidence, sources and the optional server workspace. Server observations are
labelled separately from learner-controlled local results and carry no XP or
verified track credit. Consult \path{tests/test_hosted_labs.py},
\path{tests/test_local_lab_browser.py} and
\path{scripts/verify_hosted_labs.py} for protocol and real deployment acceptance.
